\documentclass[11pt]{article}

\usepackage[margin=1in]{geometry}
\usepackage{amsmath,amssymb,amsthm,mathtools}
\usepackage{enumitem}
\usepackage{hyperref}
\hypersetup{colorlinks=true,linkcolor=blue,citecolor=blue,urlcolor=blue}
\usepackage{orcidlink}
\usepackage{fontawesome5}

\newtheorem{theorem}{Theorem}[section]
\newtheorem{proposition}[theorem]{Proposition}
\newtheorem{lemma}[theorem]{Lemma}
\newtheorem{corollary}[theorem]{Corollary}
\newtheorem{conjecture}[theorem]{Conjecture}
\newtheorem{empirical}[theorem]{Empirical Result}
\theoremstyle{remark}
\newtheorem{remark}[theorem]{Remark}
\newtheorem{example}[theorem]{Example}

\DeclareMathOperator{\ord}{ord}

\title{A Closed-Form Pressure Equation for the Accelerated\\
$qx+1$ Branching Process}

\author{Renato Augusto Tavares~{\small\href{mailto:dr.renatotavares@gmail.com}{\textcolor{black}{\faEnvelope}}}~\orcidlink{0009-0002-0196-3311}\\
\normalsize Universidade Federal de Goiás\\
\small \href{https://orcid.org/0009-0002-0196-3311}{https://orcid.org/0009-0002-0196-3311}}
\date{\today}

\begin{document}

\maketitle

\begin{abstract}
We study the reverse tree of the accelerated $qx+1$ map for odd
$q\ge3$ through a linear population functional. We prove an exact
identity, $\rho_{\mathrm{ann}}(\alpha)=q^{\alpha-1}/(2^\alpha-1)$, for
its annealed pressure, via an exact fibre-counting bijection. Besides the trivial root $\alpha=1$, the pressure equation has one
further positive root $\alpha^\ast(q)$: exactly $2$ at $q=3$, and
strictly below $1$ for every $q\ge5$. Since $\alpha^\ast(3)=2$ exceeds
$1$, the growth exponent, the smaller of the two roots, equals $1$ at
$q=3$ and equals $\alpha^\ast(q)$ for every $q\ge5$: a structural
transition at $q=5$. The $q=3$ case, exponent $1$, is the classical
Growth Exponent Conjecture of Kontorovich--Lagarias and
Applegate--Lagarias. We show that the two roots of this equation
always sit on opposite sides of a quenched/annealed freezing
transition, in the sense of directed polymers in a random environment:
the smaller root is always unfrozen, so the annealed identity
transfers to the quenched counting function of a matching i.i.d.\
branching-random-walk model; the larger root is always frozen, so the
pressure identity alone does not establish the tail index of the
density martingale. We prove that the tilted $q$-adic measure is
singular at every frozen root, by an entropy identity and a
fractional-moment estimate, and we prove the tail index for the i.i.d.\
model by an implicit renewal theorem. For the coherent $q$-adic
environment we prove martingale convergence and identify the limit
with the density of the absolutely continuous component; absolute
continuity and its tail index remain conjectural. A second, independent
conjecture states the same tail index for the growth scale factor of
the arithmetic reverse tree itself. At $q=3$ the predicted index
$\alpha^\ast=2$ matches independent empirical measurements of that
scale factor. At $q=5$, the largest of three statistical tests on the
scale factor, run at increasing sample size, gives evidence consistent
with the conjectured index, though the scale factor's own median has
not stabilized at the headroom reached; an exact test of the
martingale's own tail remains inconclusive, and neither claim is
proved for any $q\ge5$.
\end{abstract}

\section{Introduction}
\label{sec:intro}

The Collatz, or $3x+1$, conjecture generalizes naturally to a family of
$qx+1$ maps, $q$ odd. This paper studies the reverse tree of the
accelerated version of this map through a linear population functional
$G$ satisfying a natural fixed-point recursion, and derives, from an
exact bijective counting argument, a closed-form equation for the annealed growth rate of the
associated partition function. The equation's roots organize the rest
of the analysis: they govern a quenched/annealed freezing transition
familiar from directed polymers in a random environment, and their
relative position undergoes a structural change at $q=5$ that
reproduces, as a special case, a $30$-year-old open conjecture of
Kontorovich--Lagarias and Applegate--Lagarias for the classical Collatz
map.

\S\ref{sec:setup} fixes the reverse-tree construction.
\S\ref{sec:pressure} proves the pressure identity and its consequences
for the associated density martingale (absolute continuity criterion,
singularity at frozen roots, Theorem~\ref{thm:frozen-singular}).
\S\ref{sec:freezing} develops the
freezing transition. \S\ref{sec:transition} states the structural
transition at $q=5$, separating what is proved for the
matching stochastic model from what remains conjectural for the
actual arithmetic tree. \S\ref{sec:tail-regime} classifies the root
$\alpha^\ast=2$ at $q=3$ against the multivariate smoothing-transform
literature. \S\ref{sec:empirical} reports direct empirical support on
real reverse trees.

\section{Setup: the accelerated $qx+1$ map and its reverse tree}
\label{sec:setup}

Fix an odd integer $q\ge 3$ coprime to $2$. The \emph{accelerated
$qx+1$ map} is
\[
  T_q(n) = \frac{qn+1}{2^{v_2(qn+1)}}, \qquad n \text{ odd},
\]
where $v_2$ is the $2$-adic valuation. For $q=3$ this is the standard
accelerated Collatz map. The \emph{reverse tree} rooted at an odd
integer $u$ coprime to $q$ is built by, at each odd node $u$, adjoining
a child $w_a := (2^a u - 1)/q$ for every $a\ge 1$ such that
$2^a u \equiv 1 \pmod q$; such $w_a$ is automatically an odd integer
whenever it is an integer, since an odd numerator divided evenly by an
odd denominator yields an odd quotient. Writing $d:=\ord_q(2)$ for the
multiplicative order of $2$ modulo $q$, admissibility of $a$ for a
given parent $u$ depends only on $u\bmod q$: \emph{if} some
$a_0(u)\in\{1,\dots,d\}$ satisfies $2^{a_0(u)}u\equiv 1\pmod q$, it is
unique, and the admissible exponents are exactly $a\equiv a_0(u)\pmod
d$. Since $1$ is a unit mod $q$, $a_0(u)$ can exist only if $u$ itself
is a unit, so any $u$ with $\gcd(u,q)>1$, not merely $u\equiv 0\pmod
q$, is automatically \emph{sterile} (no admissible children); among
unit residues, $a_0(u)$ exists exactly when $u\bmod q$ lies in the
subgroup $\langle2\rangle\le(\mathbb{Z}/q\mathbb{Z})^\times$ generated
by $2$, and units outside $\langle2\rangle$ are sterile too. When $2$
is a primitive root mod $q$ (as for $q=3,5,9$, all used numerically
below), $\langle2\rangle$ is all of $(\mathbb{Z}/q\mathbb{Z})^\times$,
so every unit residue is fertile and the sterile classes are exactly
the non-units, i.e.\ the multiples of any prime factor of $q$: for
prime $q$ (so $q=3,5$) this is only $u\equiv0$, recovering the
familiar fact that multiples of $3$ have trivial reverse subtrees; for
$q=9=3^2$ the sterile classes are $u\equiv0,3,6\pmod9$, already beyond
the prime case. In general $d=\ord_q(2)$ can also be a proper divisor of
$\varphi(q)$ even among units: at $q=7$, $d=3<\varphi(7)=6$ and
$\langle2\rangle=\{1,2,4\}$, so $u\equiv 3,5,6\pmod7$ are sterile too,
beyond the trivial $u\equiv 0$. This does not
affect the pressure identity of Theorem~\ref{thm:pressure} below (whose
proof, via Lemma~\ref{lem:fibre-bijection}, works in the reverse
direction, from an unconstrained exponent sequence to its unique
root, and never invokes existence of $a_0(u)$ for every $u$), nor
the numerical enumeration of \S\ref{sec:empirical}, whose
implementation already excludes exactly these classes; it is a
correction to this structural description only.

We count distinct integer labels. Define
\[
 N_u(x):=\#\{n\le x:n\text{ odd and }T_q^k(n)=u
                    \text{ for some }k\ge0\}.
\]
This convention removes repetitions caused by cycles. In particular,
$N_u(x)\le\lceil x/2\rceil$.

On this tree we study the linear population functional $G$ satisfying
the fixed-point recursion
\begin{equation}\label{eq:G-recursion}
  G(u) = \sum_{a \equiv a_0(u) \ (d)} q\cdot 2^{-a}\, G(w_a),
\end{equation}
which for $q=3$ is exactly the martingale-type recursion underlying the
endogeny-barrier line of investigation developed in a companion paper
\cite{BarrierCompanion}.

\section{The pressure operator: an exact annealed identity}
\label{sec:pressure}

A first, natural attempt to organize the recursion \eqref{eq:G-recursion}
into a transfer operator is to track the parent's residue class modulo
$q$ alone. This attempt fails.

\begin{example}[The naive mod-$q$ automaton is not well defined]
\label{ex:naive-fails}
Take $q=3$ (so $d=\ord_3(2)=2$) and $a=2$, admissible for every parent
$u\equiv1\pmod3$. Two parents $u=1$ and $u=7$, both $\equiv1\pmod3$,
give $w_2(1)=(4-1)/3=1$ and $w_2(7)=(28-1)/3=9$, with $w_2(1)\equiv
1\pmod3$ but $w_2(7)\equiv0\pmod3$: the child's residue class depends
on $u\bmod9$, not merely on $u\bmod3$. In general, writing
$u=u_0+qk$ with $u_0=u\bmod q$, one finds $w_a\bmod q =
[(2^au_0-1)/q\bmod q] + [2^ak\bmod q]$: an explicit dependence on
$k=\lfloor u/q\rfloor\bmod q$. No finite modulus $q^k$ repairs this:
the same computation shows $w_a\bmod q^k$ depends on $u\bmod q^{k+1}$,
one digit further than what mod-$q^k$ knowledge of $u$ supplies. There
is therefore no well-defined finite-state transfer operator tracking a
single generation of the recursion on residue classes of any fixed
finite modulus.
\end{example}

What \emph{is} well defined, and is the correct replacement for the
naive picture of Example~\ref{ex:naive-fails}, is a bijection between
an entire admissible fibre at one level of the digit tower and the
full residue ring one level below.

\begin{lemma}[Fibre bijection]\label{lem:fibre-bijection}
Fix $a\ge1$ and $j\ge1$. The map $v\mapsto w_a(v):=(2^av-1)/q$ is a
bijection from $\{v\bmod q^j : 2^av\equiv1\!\!\pmod q\}$ (an arithmetic
progression of $q^{j-1}$ residues mod $q^j$) onto
$\mathbb{Z}/q^{j-1}\mathbb{Z}$.
\end{lemma}
\begin{proof}
Well defined: $v\bmod q^j$ determines $2^av-1\bmod q^j$, hence
$w_a(v)\bmod q^{j-1}$. Injective: if $v\not\equiv v'\pmod{q^j}$ within
the progression, write $v-v'=qt$ with $t\not\equiv0\pmod{q^{j-1}}$;
then $w_a(v)-w_a(v')=2^at\not\equiv0\pmod{q^{j-1}}$ since $2$ is a unit
mod $q^{j-1}$. Surjectivity follows since the map is an injection
between two finite sets of equal size $q^{j-1}$.
\end{proof}

Iterating Lemma~\ref{lem:fibre-bijection} along a depth-$k$ path
$(a_1,\dots,a_k)$ shows every sequence $(a_1,\dots,a_k)\in
(\mathbb{Z}_{\ge1})^k$, with no admissibility restriction on the
individual $a_i$, is admissible for \emph{exactly one} root residue
$u_0\bmod q^k$, namely
$u_0\equiv\sum_{i=1}^k q^{i-1}2^{-S_i}\pmod{q^k}$, $S_i:=a_1+\cdots+a_i$.
Define the depth-$k$ partition function rooted at $u_0$,
\[
  Z_k(\alpha;u_0) := \sum_{\substack{(a_1,\dots,a_k)\text{ admissible}\\ \text{from } u_0}} \ \prod_{i=1}^k \bigl(q\cdot2^{-a_i}\bigr)^\alpha .
\]

\begin{theorem}[Exact annealed pressure identity]\label{thm:pressure}
For every odd $q\ge3$, every $k\ge1$, and every $\alpha>0$,
\begin{equation}\label{eq:annealed-identity}
  \sum_{u_0\in\mathbb{Z}/q^k\mathbb{Z}} Z_k(\alpha;u_0) \;=\; \Bigl(\frac{q^\alpha}{2^\alpha-1}\Bigr)^{k},
\end{equation}
so the \emph{average} of $Z_k(\alpha;\cdot)$ over all $q^k$ root
residues equals exactly $\bigl(q^{\alpha-1}/(2^\alpha-1)\bigr)^k$; the
associated annealed pressure is
\begin{equation}\label{eq:pressure-closed-form}
  \rho_{\mathrm{ann}}(\alpha) = \frac{q^{\alpha-1}}{2^\alpha-1},
\end{equation}
giving the pressure equation at $\rho_{\mathrm{ann}}=1$:
\begin{equation}\label{eq:pressure-eq}
  q^{\alpha-1} = 2^{\alpha}-1.
\end{equation}
\end{theorem}
\begin{proof}
By Lemma~\ref{lem:fibre-bijection} (iterated), the admissible depth-$k$
sequences rooted at the $q^k$ distinct residues $u_0$ partition the
full, unconstrained set $(\mathbb{Z}_{\ge1})^k$ with no overlap and no
omission. Summing $\prod_i(q2^{-a_i})^\alpha$ over \emph{all} sequences
in $(\mathbb{Z}_{\ge1})^k$ therefore factors as
$\bigl(\sum_{a=1}^\infty (q2^{-a})^\alpha\bigr)^k =
\bigl(q^\alpha/(2^\alpha-1)\bigr)^k$, which is exactly the left-hand
side of \eqref{eq:annealed-identity} summed over the partition.
\end{proof}

We verified \eqref{eq:annealed-identity} independently by direct
enumeration for $q\in\{3,5,7,9\}$, $k\in\{1,2,3\}$, and
$\alpha\in\{0.26,0.5,1,1.37,2\}$, matching to the truncation error of
the (rapidly convergent) tail of the sum over $a$.

\begin{remark}[Mass conservation as a corollary]
At $\alpha=1$, the annealed pressure \eqref{eq:pressure-closed-form}
evaluates directly to $\rho_{\mathrm{ann}}(1)=q^0/(2^1-1)=1$ for every
$q$ (mass conservation of the branching process). Thus $\alpha=1$ is
always a root of \eqref{eq:pressure-eq}.
\end{remark}

\begin{lemma}[Log-convexity of the pressure function]\label{lem:log-convex}
For every odd $q\ge3$, $P(\alpha):=\log\rho_{\mathrm{ann}}(\alpha)$ is
strictly convex on $\alpha>0$. Consequently
equation~\eqref{eq:pressure-eq} has at most two positive roots, and if
it has exactly two, $\alpha_-<\alpha_+$, then $P'(\alpha_-)<0<P'(\alpha_+)$.
\end{lemma}
\begin{proof}
From \eqref{eq:pressure-closed-form}, $P(\alpha)=(\alpha-1)\log q-\log(2^\alpha-1)$, so
\[
 P''(\alpha)=\frac{\log^2 2\cdot 2^\alpha}{(2^\alpha-1)^2}>0
\]
for every $\alpha>0$: strict convexity. A strictly convex function
equals any given value at most twice, and, at two such roots, must be
decreasing through the first and increasing through the second (on
either side of its unique minimum), giving $P'(\alpha_-)<0<P'(\alpha_+)$.
\end{proof}

At a pressure root $\alpha$, the numbers $p_\alpha(a):=q^{\alpha-1}2^{-\alpha a}$,
$a\ge1$, sum to one, so they define a probability law on branch
exponents. For independent digits $a_1,a_2,\ldots$ drawn from
$p_\alpha$, the partial sums
$F_k:=\sum_{i=1}^k q^{i-1}2^{-S_i}\pmod{q^k}$, $S_i:=a_1+\cdots+a_i$,
form a projectively consistent family (Lemma~\ref{lem:fibre-bijection}
gives $F_{k+1}\equiv F_k\pmod{q^k}$), hence converge to a
$\mathbb{Z}_q$-valued limit; write $\mu_\alpha$ for its law, the
\emph{tilted Syracuse measure} at $\alpha$.

\begin{theorem}[$q$-adic density martingale]\label{thm:qadic-martingale}
Let $\alpha$ be either positive root of \eqref{eq:pressure-eq}, and let
$U$ be Haar-uniform on $\mathbb Z_q$.
Then
\[
 M_k^{(\alpha)}(U):=Z_k(\alpha;U\bmod q^k)
\]
is a nonnegative martingale of mean one for the filtration generated by
$U\bmod q^k$. Consequently $M_k$ converges almost surely to a finite
limit $W_{q,\alpha}$. Writing
$\mu_\alpha=f_{q,\alpha}\,m_q+\mu_{\mathrm s}$ for the Lebesgue
decomposition of the tilted Syracuse measure above
with respect to Haar measure $m_q$, then
\[
 W_{q,\alpha}=f_{q,\alpha}(U)\qquad m_q\text{-almost surely}.
\]
In particular, $W_{q,\alpha}$ is nonzero with positive probability
exactly when $\mu_\alpha$ has a nonzero absolutely continuous component,
and $\mathbb E W_{q,\alpha}=1$ exactly when $\mu_\alpha\ll m_q$.
\end{theorem}

\begin{proof}
With $F_k$ and $p_\alpha$ as above, the fibre identity gives
\[
 Z_k(\alpha;u)=q^k\Pr(F_k=u).
\]
Conditional on $U\equiv u\pmod {q^k}$,
\[
 \mathbb E[M_{k+1}\mid U\bmod q^k=u]
 =\frac1q\sum_{v\equiv u\ (q^k)}q^{k+1}\Pr(F_{k+1}=v)
 =q^k\Pr(F_k=u)=M_k(u).
\]
The nonnegative martingale convergence theorem gives the limit. The
functions $M_k$ are the Radon--Nikodym derivatives of the restrictions
of $\mu_\alpha$ to the cylinder sigma-algebras relative to the
corresponding restrictions of $m_q$. Martingale differentiation of the
Lebesgue decomposition identifies their almost-sure limit with
$d(\mu_\alpha)_{\mathrm{ac}}/dm_q$. Its integral is the total mass of
the absolutely continuous part, which proves the final assertions.
\end{proof}

\begin{theorem}[$L^p$ collision criterion]\label{thm:lp-collision}
Let $\alpha$ be a pressure root, let $p>1$, let $\mu_{\alpha,k}$ be the
law of $F_k$ modulo $q^k$, and
let $M_k^{(\alpha)}=q^k\mu_{\alpha,k}$ be its density relative to
uniform measure on $\mathbb Z/q^k\mathbb Z$. Then
\begin{equation}\label{eq:lp-collision}
 \lVert M_k^{(\alpha)}\rVert_p^p
 =q^{k(p-1)}\sum_{u\bmod q^k}\mu_{\alpha,k}(u)^p.
\end{equation}
The projective measure $\mu_\alpha$ has a density in $L^p(\mathbb Z_q)$
if and only if the right-hand side is bounded uniformly in $k$. If $p$
is an integer, the sum in \eqref{eq:lp-collision} is the probability
that $p$ independent tilted Syracuse sums agree modulo $q^k$.
\end{theorem}

\begin{proof}
Uniform measure assigns mass $q^{-k}$ to each residue, so
\[
 \lVert M_k^{(\alpha)}\rVert_p^p
 =q^{-k}\sum_u\bigl(q^k\mu_{\alpha,k}(u)\bigr)^p,
\]
which is \eqref{eq:lp-collision}. For integer $p$, expanding the
probability of equality of $p$ independent copies gives the collision
interpretation.

If the quantities in \eqref{eq:lp-collision} are uniformly bounded,
the martingale $M_k^{(\alpha)}$ converges in $L^p$. Its limit has integral
one and is the density of the projective measure on every cylinder, so
$\mu_\alpha$ is absolutely continuous with an $L^p$ density. Conversely,
if $d\mu_\alpha=f\,dm_q$ with $f\in L^p$, then
$M_k^{(\alpha)}=\mathbb E[f\mid U\bmod q^k]$. Conditional expectation is
an $L^p$ contraction, which gives the uniform bound.
\end{proof}

For $q=3$ and $\alpha=1$, Theorem~\ref{thm:lp-collision} at $p=2$
gives exactly the collision statistic used in the finite-level $L^2$
growth measurement of the endogeny-barrier companion paper
\cite{BarrierCompanion}. Thus that $L^2$ sufficient condition is one
member of a complete family of weighted collision criteria
(finite-level spectrum computed in the accompanying repository). These
are averaged concentration conditions; they do not control the
smallest cylinder mass required by a worst-residue covering problem.

\begin{theorem}[Singularity at the frozen root]
\label{thm:frozen-singular}
Let $\alpha=\alpha_+(q)$, the larger of the two positive roots of
\eqref{eq:pressure-eq} (\S\ref{sec:freezing}). Then
\[
 M_k^{(\alpha)}(U)\longrightarrow0
 \qquad m_q\text{-almost surely},
\]
and the tilted Syracuse measure $\mu_\alpha$ is singular with respect
to Haar measure.
\end{theorem}

\begin{proof}
Write $r=2^{-\alpha}$. The digit law is
$p_\alpha(a)=(1-r)r^{a-1}$, so its entropy is
\[
 \mathrm{Ent}(p_\alpha)=-\log(1-r)-\frac{r}{1-r}\log r.
\]
At a pressure root, direct differentiation gives
\[
 \mathrm{Ent}(p_\alpha)-\log q
 =-\alpha P'(\alpha)=s(\alpha).
\]
The larger root is frozen, hence $\mathrm{Ent}(p_\alpha)<\log q$ (see
\S\ref{sec:freezing}). For $0<t<1$,
subadditivity of $x^t$ and the representation
$M_k^{(\alpha)}(u)=q^k\sum_{F_k=u}\prod_i p_\alpha(a_i)$ give
\[
 \mathbb E_{m_q}\bigl[(M_k^{(\alpha)})^t\bigr]
 \le\left(q^{t-1}\sum_{a\ge1}p_\alpha(a)^t\right)^k.
\]
The logarithmic derivative of the factor in parentheses at $t=1$ is
$\log q-\mathrm{Ent}(p_\alpha)>0$. It is therefore less than one for some
$t<1$ sufficiently close to one. The fractional moment tends to zero,
so the almost-sure martingale limit is zero. Theorem~\ref{thm:qadic-martingale}
then identifies the absolutely continuous component as zero.
\end{proof}

\section{Quenched vs.\ annealed: a freezing transition}
\label{sec:freezing}

Theorem~\ref{thm:pressure} is a statement about the \emph{average} of
$Z_k(\alpha;\cdot)$ over root residues; it is not, by itself, a
statement about $Z_k(\alpha;u_0)^{1/k}$ for a single, fixed root
$u_0$, which is what actually governs the population martingale
$G(v)$ for one specific integer $v$. Whether the two coincide is
exactly the quenched-versus-annealed question familiar from directed
polymers in a random environment and from Derrida's random energy
model, governed by a standard freezing criterion. Writing
$P(\alpha):=\log\rho_{\mathrm{ann}}(\alpha)$ for the annealed
log-pressure and $s(\alpha):=P(\alpha)-\alpha P'(\alpha)$ for the
associated entropy, quenched and annealed rates coincide where
$s(\alpha)>0$ (an \emph{unfrozen} phase) and diverge where
$s(\alpha)<0$ (a \emph{frozen} phase, in which rare, low-entropy paths
dominate the quenched rate instead of the bulk of the ensemble).

\begin{proposition}[The larger root is always frozen]\label{prop:always-frozen}
For every odd $q\ge3$, writing $\alpha_-<\alpha_+$ for the two positive
roots of $\rho_{\mathrm{ann}}(\alpha)=1$: $s(\alpha_-)>0$ (unfrozen) and
$s(\alpha_+)<0$ (frozen), unconditionally.
\end{proposition}
\begin{proof}
At any root, $P(\alpha)=0$, so $s(\alpha)=-\alpha P'(\alpha)$.
Lemma~\ref{lem:log-convex} gives $P'(\alpha_-)<0<P'(\alpha_+)$
directly. Hence
$s(\alpha_-)=-\alpha_-P'(\alpha_-)>0$ and
$s(\alpha_+)=-\alpha_+P'(\alpha_+)<0$.
\end{proof}

For $q\ge5$, where $\alpha_+=1$ always, this has an exact closed form:
at $\alpha=1$, $2^\alpha-1=1$ so $\log(2^\alpha-1)=0$, and
$s(1)=2\log2-\log q=\log(4/q)$, which is negative exactly when $q>4$.
So the freezing of the trivial root for every odd $q\ge5$ is not a
numerical coincidence but the same sign condition as the mean-drift
remark below ($\log q-2\log2>0$): the drift that makes orbits expand on
average is, up to a minus sign, exactly the entropy deficit that
freezes $\alpha_+=1$.

Since $s'(\alpha)=P'(\alpha)-P'(\alpha)-\alpha P''(\alpha)=-\alpha
P''(\alpha)<0$ for every $\alpha>0$ by Lemma~\ref{lem:log-convex}, $s$
is strictly decreasing throughout, so its root $\alpha_c(q)$ is unique
wherever it exists; existence, and that it lies strictly between
$\alpha_-(q)$ and $\alpha_+(q)$, follows from
Proposition~\ref{prop:always-frozen} by the intermediate value theorem.
We computed $\alpha_c(q)$ directly:
\begin{center}
\begin{tabular}{c|c}
$q$ & $\alpha_c(q)$ \\ \hline
$3$ & $1.355$ \\
$5$ & $0.794$ \\
$7$ & $0.564$ \\
$9$ & $0.437$
\end{tabular}
\end{center}
confirming $\alpha_-(q)<\alpha_c(q)<\alpha_+(q)$ in every case, as
Proposition~\ref{prop:always-frozen} requires. We also verified
numerically, with an independent memoised recursion computing
$Z_k(\alpha;u_0)$ exactly at a fixed (not averaged) root, that the
freezing reflects a real divergence of quenched from annealed growth:
at $q=3$, $u_0=1$ (a period-one cycle of the residue dynamics, since
$w_2(1)=1$ exactly), $Z_k(2;1)^{1/k}$ shows an oscillation with
geometric-mean ratio well below the annealed rate
$\rho_{\mathrm{ann}}(2)=1$; at $q=5$, $Z_k(1;u_0)^{1/k}$ at the roots
$u_0\in\{1,2,3,4\}$ likewise fails to track $\rho_{\mathrm{ann}}(1)=1$.
This check has a limitation: every one of these
roots lies on or feeds directly into the smallest-admissible-branch
map's own trivial cycle ($1\leftrightarrow3$ for $q=5$, exactly as
$u_0=1$ does for $q=3$), so the divergence observed here has a
deterministic-recurrence explanation available in addition to the
entropy mechanism of Proposition~\ref{prop:always-frozen}, and does
not by itself sample the generic (non-periodic) environment that the
quenched statement is really about. This illustrates freezing at a
special, easily-computed class of roots. \S\ref{sec:empirical}
below carries the corresponding generic-environment evidence, a
handful of roots at $q=5,7$ sampled well away from any known short
cycle.

The pressure identity transfers to the counting exponent of the i.i.d.
model at the unfrozen root $\alpha_-(q)$. It does not by itself prove a
tail law. Theorem~\ref{thm:iid-tail} obtains the i.i.d. tail by implicit
renewal, while Conjecture~\ref{conj:tail-index} states the separate
claim for a coherent Haar-random $q$-adic root environment.

\section{The structural transition at $q=5$}
\label{sec:transition}

Besides the trivial root $\alpha=1$, equation~\eqref{eq:pressure-eq}
has exactly one further positive root $\alpha^\ast(q)$. By
Lemma~\ref{lem:log-convex}, $\rho_{\mathrm{ann}}(\alpha)=1$ has at
most two positive roots; for existence, $\rho_{\mathrm{ann}}(\alpha)\to
\infty$ both as $\alpha\to0^+$ (since $2^\alpha-1\to0$ linearly while
$q^{\alpha-1}\to1/q$) and as $\alpha\to\infty$ (since $q>2$), so
$P:=\log\rho_{\mathrm{ann}}\to\infty$ at both ends, and
$P'(1)=\log(q/4)\ne0$ for every odd $q$ (\S\ref{sec:freezing}) places
the strictly convex function's unique minimum away from $\alpha=1$,
forcing a second root on the side of $1$ that the sign of $P'(1)$
selects. We solved for $\alpha^\ast(q)$ independently by bisection:

\begin{center}
\begin{tabular}{c|c|c}
$q$ & second root $\alpha^\ast(q)$ & interpretation \\ \hline
$3$ & $2.000000$ & exact \\
$5$ & $0.650919$ & \\
$7$ & $0.373501$ & \\
$9$ & $0.258108$ &
\end{tabular}
\end{center}

At $q=3$ the second root is exactly $\alpha^\ast=2$, matching the
universal tail exponent for the classical Collatz reverse tree,
established directly on real reverse-tree samples by independent
methods (a Hill estimator and extreme value theory on block maxima).
For $q\ge 5$, however, the second root falls
strictly \emph{below} $1$, a qualitative structural transition. Writing
$\alpha_-<\alpha_+$ for
the two roots, we have $\{\alpha_-,\alpha_+\}=\{1,2\}$ at $q=3$ but
$\{\alpha_-,\alpha_+\}=\{\alpha^\ast(q),1\}$ for $q\ge5$: the trivial
root becomes the \emph{larger} of the two.

\paragraph{The i.i.d.\ model.} Idealize the reverse tree by declaring
each candidate exponent $a\ge1$ independently fertile with probability
$1/q$; a fertile $a$ produces a child at multiplicative displacement
$2^a/q$ from its parent, matching the arithmetic tree's own
branch-by-branch transition rule (\S\ref{sec:setup}). Writing
$S_v\in\mathbb{R}$ for the cumulative log-displacement
of an individual $v$ from the root, and $N(x):=\#\{v:S_v\le\log x\}$ for
the number of individuals of any generation reaching size at most $x$,
the many-to-one formula gives, for $\theta>0$,
\[
\mathbb E\Big[\sum_{|v|=n}e^{-\theta S_v}\Big]
=\Big(\mathbb E\Big[\sum_{a\ge1}\mathbf1[\text{fertile}]\,(q/2^a)^{\theta}\Big]\Big)^{n}
=\Big(\tfrac1q\sum_{a\ge1}(q/2^a)^\theta\Big)^{n}
=\rho_{\mathrm{ann}}(\theta)^n,
\]
so the model's per-generation growth rate is governed by exactly the
same pressure function as Theorem~\ref{thm:pressure}.

\begin{theorem}[Structural transition: exponent, i.i.d.\ model]\label{thm:transition-model}
For the i.i.d.\ model above, $\log N(x)/\log x\to\alpha_-(q)$ almost
surely, for every odd $q\ge3$: exponent $1$ (i.e.\ $N(x)=x^{1+o(1)}$)
at $q=3$; exponent $\alpha_-(q)<1$ for every $q\ge5$.
\end{theorem}

\begin{proof}
\emph{Upper bound.} Fix $\theta\in(\alpha_-(q),\alpha_+(q))$, where
$\rho_{\mathrm{ann}}(\theta)<1$ by the log-convexity argument above (a
convex function vanishing at two points $\alpha_-<\alpha_+$ is
non-positive in between). By Markov's inequality applied to the
many-to-one identity, $\mathbb E[\#\{|v|=n:S_v\le y\}]\le
e^{\theta y}\rho_{\mathrm{ann}}(\theta)^n$, so summing over all
generations, $\mathbb E[N(e^y)]\le e^{\theta y}/(1-\rho_{\mathrm{ann}}(\theta))$.
A second Markov step gives $\Pr(N(x)>x^{\theta+\varepsilon})\le
x^{-\varepsilon}/(1-\rho_{\mathrm{ann}}(\theta))$ for any
$\varepsilon>0$; this is summable along $x=2^k$, so Borel--Cantelli and
monotonicity of $N$ give $\limsup_x\log N(x)/\log x\le\theta+\varepsilon$
almost surely. Letting $\theta\downarrow\alpha_-(q)$ and
$\varepsilon\downarrow0$ gives $\limsup\le\alpha_-(q)$.

\emph{Lower bound.} Write $\alpha:=\alpha_-(q)$ and
$\mu:=-\rho_{\mathrm{ann}}'(\alpha)>0$ (unfrozen,
Proposition~\ref{prop:always-frozen}). Let $a_1,a_2,\ldots$ be i.i.d.\
draws from the digit law $p_\alpha(a):=q^{\alpha-1}2^{-\alpha a}$ of
\S\ref{sec:pressure}, and set $\xi_i:=a_i\log2-\log q$, the induced
displacement per step. Since $p_\alpha$ is geometric with mean
$\mathbb E[a_i]=2^\alpha/(2^\alpha-1)$,
\[
 \mathbb E[\xi_i]=\log2\cdot\frac{2^\alpha}{2^\alpha-1}-\log q
 =-P'(\alpha)=\mu,
\]
so the $\xi_i$ are i.i.d.\ with mean exactly $\mu$. Fix $0<\varepsilon<\mu$
and $L\ge1$, and let $Y_L$ be the
number of generation-$L$ descendants of a node whose cumulative
displacement lies in $[L(\mu-\varepsilon),L(\mu+\varepsilon)]$. The
many-to-one identity at $\theta=\alpha$ (where $\rho_{\mathrm{ann}}(\alpha)=1$)
gives
\[
 r_L:=\mathbb E[Y_L]
 =\mathbb E_\alpha\Bigl[e^{\alpha(\xi_1+\cdots+\xi_L)}
   \mathbf 1_{\{\xi_1+\cdots+\xi_L\in[L(\mu-\varepsilon),L(\mu+\varepsilon)]\}}\Bigr]
 \ge e^{\alpha L(\mu-\varepsilon)}
   \Pr_\alpha\Bigl(\tfrac1L\textstyle\sum_i\xi_i\in[\mu-\varepsilon,\mu+\varepsilon]\Bigr),
\]
and the law of large numbers sends the probability factor to $1$ as
$L\to\infty$, so $r_L\to\infty$; fix $L$ with $r_L>1$. For any $K$,
keep only the first $K$ such descendants at every node (in some fixed
measurable order): this defines a sub-population, resampled
every $L$ generations, that is an ordinary Galton--Watson process with
bounded offspring $Y_L^{(K)}:=\min(Y_L,K)\in\{0,\ldots,K\}$, hence of
finite variance, and mean $r_L^{(K)}:=\mathbb E[Y_L^{(K)}]\uparrow r_L$
as $K\to\infty$ by monotone convergence; fix $K$ with $r_L^{(K)}>1$.
The classical growth theorem for a finite-variance supercritical
Galton--Watson process gives, on its (positive-probability) survival
event, $\frac1j\log Z_j^{(K)}\to\log r_L^{(K)}$ almost surely, where
$Z_j^{(K)}$ is the population after $j$ such blocks. Since the root has
infinitely many independent admissible children almost surely (the
events $2^au\equiv1\!\!\pmod q$, $a\ge1$, are i.i.d.\ with probability
$1/q$ each), and each spawns an independent copy of the whole
construction, at least one copy survives almost surely, rooted at an
almost surely finite generation and displacement
$\Delta<\infty$. Every one of the $Z_j^{(K)}$ retained individuals then
lies at generation $jL+O(1)$ with displacement at most
$jL(\mu+\varepsilon)+\Delta$, so
$N\bigl(e^{jL(\mu+\varepsilon)+\Delta}\bigr)\ge Z_j^{(K)}$, and
monotonicity of $N$ (the exponentially spaced points
$jL(\mu+\varepsilon)+\Delta$ have consecutive ratio $\to1$, since
$\Delta$ is fixed while $j\to\infty$) gives
\[
 \liminf_{x\to\infty}\frac{\log N(x)}{\log x}
 \ge\frac{\log r_L^{(K)}}{L(\mu+\varepsilon)}\qquad\text{almost surely.}
\]
This holds simultaneously for every $K$ (a countable intersection of
almost-sure events, one per $K\in\mathbb N$), so the right side may be
replaced by its supremum over $K$, namely $\log r_L/(L(\mu+\varepsilon))$;
call this almost-sure event $\mathcal E_L$. Intersecting
$\mathcal E_L$ over the countably many $L\in\mathbb N$ with $r_L>1$
gives an almost-sure event on which the bound holds simultaneously for
every such $L$, so $L\to\infty$ may be taken along this fixed sequence
of events, followed by $\varepsilon\downarrow0$; using
$\log r_L\ge\alpha L(\mu-\varepsilon)+\log\Pr_\alpha(\cdots)\to\alpha
L(\mu-\varepsilon)(1+o(1))$, this gives
$\liminf_x\log N(x)/\log x\ge\alpha$ almost surely. With the upper
bound, $\log N(x)/\log x\to\alpha_-(q)$ almost surely.
\end{proof}

\begin{remark}[Relation to Kontorovich--Lagarias's Theorem~8.10]
Kontorovich--Lagarias \cite[Theorem~8.10]{KontorovichLagarias2010}
prove an analogous exact limit, for $q=5$, for a different
branching-random-walk representation $B[5^0]$ of the $5x+1$ reverse
tree, one with at most two offspring per node rather than the present
model's almost-surely infinite offspring; a direct computation shows
its log-Laplace transform is $m_{\mathrm{KL}}(\theta)=
2^{-\theta}(1+5^{\theta-1})$, not $\rho_{\mathrm{ann}}(\theta)$. They
share every root: $m_{\mathrm{KL}}(\theta)-1=\bigl((2^\theta-1)/2^\theta\bigr)
\bigl(\rho_{\mathrm{ann}}(\theta)-1\bigr)$, which is why
$\eta_{5,BP}=\alpha_-(5)$ in Remark~\ref{rem:novelty-109} despite the
two derivations having nothing else in common; their proof does not
transfer to the present model, and is not used above.
\end{remark}

\begin{conjecture}[Finer asymptotic for the counting function]\label{conj:transition-fine}
The counting function satisfies the sharper asymptotic
$N(x)\sim W\cdot x^{\alpha_-(q)}$, not merely
Theorem~\ref{thm:transition-model}'s exponent statement, holding almost
surely and uniformly, \emph{including} at $q=3$, for an almost surely
positive, non-degenerate random scale factor $W$ (not a deterministic
constant at any $q$, since the underlying reproduction law has nonzero
variance). This is a strictly stronger claim than
Theorem~\ref{thm:transition-model}: it pins down the constant in front
of $x^{\alpha_-(q)}$, which needs machinery
beyond the exponential tilting and block truncation used above. A directed literature
search found the neighboring machinery but not this exact statement
ready-made: the classical Crump--Mode--Jagers renewal theorem
\cite{Nerman1981} is stated for one-sided ($[0,\infty)$) reproduction
point processes, while our branching displacements include negative
values (e.g.\ $2/3<1$ at $q=3$, since a child can be numerically
smaller than its parent); the branching-random-walk martingale theory
that is natively two-sided \cite{Biggins1992} gives convergence
generation-by-generation, not a renewal law summed over all generations;
and the implicit renewal theory for weighted trees with general
(two-sided) weights \cite{Goldie1991,JelenkovicOlveraCravioto2012,
AlsmeyerBigginsMeiners2012} resolves the \emph{tail} of the fixed-point
equation for $W$ itself, not the cumulative counting law for $N(x)$, a
related but distinct question. We expect the present statement to
follow from combining this machinery (spine decomposition and
exponential tilting for a two-sided one-dimensional random walk,
exactly as already used for the tail-of-$W$ results just cited) but we
have not carried out that adaptation, nor located it already written in
this exact form. A very recent preprint \cite{VillemonaisZalduendo2025}
extends Biggins-type martingale convergence to general type spaces and
non-geometric weights, but at the level of individual generations, the
same scope as \cite{Biggins1992} above; it contains no renewal or
Crump--Mode--Jagers machinery and does not address the cumulative,
summed-over-generations statement needed here. The non-lattice
hypothesis that any such extension of
\cite{Nerman1981} would need holds here, and it is exactly the
irrationality of $\log q/\log2$ that secures it. Although each
single-step displacement $-\log q + a\log2$ ($a\in\mathbb{Z}_{\ge1}$)
lies in one coset of the lattice $\log(2)\,\mathbb{Z}$, containment in
a coset is not arithmeticity: the smallest closed subgroup of
$\mathbb{R}$ generated by these displacements is
$\log(q)\,\mathbb{Z}+\log(2)\,\mathbb{Z}$, which is dense in
$\mathbb{R}$ for every odd $q\ge3$ because $\log q/\log2\notin\mathbb{Q}$.
Equivalently, the characteristic function of the displacement equals
$1$ only at the origin, so the renewal/CMJ asymptotic carries no
log-periodic (period $\log2$) modulation. Were $\log q/\log2$ rational
the branch positions would collapse onto a lattice and a
periodic correction would appear.
\end{conjecture}

\begin{conjecture}[Structural transition: exponent, arithmetic tree]\label{conj:transition-arithmetic}
For every admissible root $u$, the actual reverse tree of $T_q$
satisfies
\[
 N_u(x)=x^{\alpha_-(q)+o(1)}.
\]
Thus the exponent is $1$ for $q=3$ and is strictly below $1$ for every
odd $q\ge5$. At $q=3$ this is the Growth Exponent Conjecture of
Kontorovich--Lagarias \cite{KontorovichLagarias2010} and
Applegate--Lagarias
\cite{ApplegateLagarias1995I,ApplegateLagarias1995II} for the classical
Collatz reverse tree, namely $\pi_a(x)=x^{1+o(1)}$ for their analogue
of $N_u(x)$ above (the count, up to $x$, of odd integers reaching a
fixed root $a$ under repeated reverse iteration), open
since 1995, with the best unconditional partial result remaining the
one-sided bound $\pi_a(x)\ge x^{0.84}$ of Krasikov--Lagarias
\cite{KrasikovLagarias2003}. Our pressure identity reproduces the
exponent-$1$ statement of this conjecture as the $\alpha_-=1$ case of
a single closed-form equation, together with its $q\ge5$ analogue; it
does not resolve the conjecture. A calibrated empirical test of this
conjecture at $q=5$ against a competing prediction is the subject of a
companion paper \cite{KLVolkovCompanion}.
\end{conjecture}

We state Theorem~\ref{thm:transition-model},
Conjecture~\ref{conj:transition-fine}, and
Conjecture~\ref{conj:transition-arithmetic} as three separate claims,
not one theorem about the arithmetic tree, because each step up in
strength, from the proved exponent to the conjectural finer asymptotic
to the conjectural transfer to the arithmetic tree itself, is exactly
the kind of transfer that
Example~\ref{ex:naive-fails} shows can fail;
Remark~\ref{rem:transfer-basis} below specifies what does and does
not rest on proof. Direct empirical support for
Conjecture~\ref{conj:transition-arithmetic} is given in
\S\ref{sec:empirical} below.

\begin{remark}[What ``transfers'' rests on]\label{rem:transfer-basis}
For the i.i.d.\ branching-random-walk model that the annealed
identity of Theorem~\ref{thm:pressure} describes, quenched-equals-annealed
growth throughout the unfrozen region is a rigorous, classical result
\cite{Biggins1992}. For our arithmetic recursion, no analogous theorem
is known to us: the closest arithmetic result in
the literature is a one-sided density bound of Applegate--Lagarias
\cite{ApplegateLagarias1995I,ApplegateLagarias1995II}, later sharpened
by Krasikov--Lagarias \cite{KrasikovLagarias2003}, not a matching
two-sided quenched statement. We treat the transfer, at every root we
invoke it for, as supported by (i) the classical stochastic-model
theorem, (ii) direct numerical confirmation, and not by a proof for the
arithmetic tree itself.
\end{remark}

\begin{conjecture}[Density and tail of the tilted Syracuse measure]
\label{conj:tail-index}
Let $U$ be Haar-uniform on $\mathbb Z_q$, and read
$Z_k(\alpha_-(q);U)$ through its residue modulo $q^k$. The sequence
has the limit $W_q:=W_{q,\alpha_-(q)}$ of
Theorem~\ref{thm:qadic-martingale}. The tilted
measure is absolutely continuous, its density $W_q$ is non-degenerate,
and
\[
 \Pr(W_q>x)=x^{-\alpha_+(q)/\alpha_-(q)+o(1)}.
\]
Here $\alpha_+(q)$ is the larger root of \eqref{eq:pressure-eq}
($\alpha_+(3)=2$; $\alpha_+(q)=1$ for every $q\ge5$).
\end{conjecture}

\begin{theorem}[Tail index in the i.i.d. model]\label{thm:iid-tail}
For the additive martingale at parameter $\alpha_-(q)$ in the matching
i.i.d. branching model,
\[
 \Pr(W_\infty>x)\sim C_q x^{-\alpha_+(q)/\alpha_-(q)}
 \qquad (x\to\infty)
\]
for a constant $C_q>0$.
\end{theorem}

\begin{proof}
Apply the implicit renewal theorem for generalized multiplicative
cascades \cite{Liu2000}. With $\theta=\alpha_-(q)$, its logarithmic
moment function is
\[
 \psi(\sigma)=\log\rho_{\mathrm{ann}}(\theta \sigma).
\]
It vanishes at $\sigma=1$ and next at
$\kappa=\alpha_+(q)/\alpha_-(q)$, while $\psi'(1)<0$ because the smaller
root is unfrozen. The geometric offspring ladders give all moments
needed in a neighborhood of $\kappa$. The closed subgroup generated by
the logarithmic weights is generated by $\theta\log2$ and
$\theta\log q$; it is dense because $\log q/\log2$ is irrational.
Thus the renewal law is non-arithmetic and has a constant, positive
tail coefficient.
\end{proof}

The hypotheses of Liu's theorem invoked above, namely
$\rho_{\mathrm{ann}}(\alpha_-(q))=1$, $\rho_{\mathrm{ann}}(\alpha_+(q))=1$,
$\psi'(1)<0$, and $\kappa=\alpha_+(q)/\alpha_-(q)>1$, were checked
numerically for $q=3,5,7,9,15$ in the accompanying repository.

This is the classical smoothing-transform result for the stochastic
model. Its transfer to the coherent $q$-adic root environment is
Conjecture~\ref{conj:tail-index}, since the arithmetic subtrees are not
independent.
The tail-index equation $\mathbb{E}[\sum_i A_i^\sigma]=1$ becomes,
for our weights, $q^{\alpha_-\sigma-1}=2^{\alpha_- \sigma}-1$, solved by
$\sigma=1$ trivially and, nontrivially, by
$\sigma=\alpha_+/\alpha_-$. This is further reason to take the conjecture
seriously as a prediction, even though, per
Remark~\ref{rem:transfer-basis}, the theory rigorously covers only the
stochastic model so far.

We state this as a conjecture separate from both
Theorem~\ref{thm:transition-model} and
Conjecture~\ref{conj:transition-arithmetic} above, because
Proposition~\ref{prop:always-frozen} shows $\alpha_+(q)$ is
\emph{always} frozen, so the annealed pressure identity alone does not
establish the tail. Theorem~\ref{thm:iid-tail} obtains it by implicit
renewal at the unfrozen martingale normalization. Its arithmetic
counterpart still requires evidence external to
Theorem~\ref{thm:pressure}.

Every numerical test below samples actual integer roots of the
arithmetic tree, not points of the Haar-uniform $q$-adic population
$U$ that Conjecture~\ref{conj:tail-index} is stated over. What these
tests bear on directly is a different, equally open claim: the tail of
the scale factor already introduced in
Conjecture~\ref{conj:transition-fine}.

\begin{conjecture}[Tail of the real-tree growth scale factor]\label{conj:real-tree-tail}
Fix odd $q\ge3$ and headroom $H\ge1$. Let $v$ range uniformly over the
admissible roots in an interval growing to infinity, and set
$W_v(H):=N_v(vH)/H^{\alpha_-(q)}$ ($N_v$ as in \S\ref{sec:setup}); as
the interval grows, $W_v(H)$ converges in law to a limit $W(H)$. As
$H\to\infty$, $W(H)$ converges in law to the scale factor $W$ of
Conjecture~\ref{conj:transition-fine}, with upper tail
\[
 \Pr(W>x) = x^{-\alpha_+(q)/\alpha_-(q)+o(1)}.
\]
\end{conjecture}

The same implicit-renewal mechanism that proves Theorem~\ref{thm:iid-tail}
for the i.i.d.\ model predicts this exponent for $W$,\footnote{A related
structural prediction of the multitype model,
confirmed empirically for $q=5$: conditioning on the root's residue
class $i=u_0\bmod q$ (with $a_0(i)$ the smallest admissible branch
exponent for class $i$), $W$ decomposes by class, at the tail exponent
$\theta:=\alpha_-(q)$ of Conjecture~\ref{conj:real-tree-tail}, as
$W_i\stackrel{d}{=}2^{-a_0(i)\theta}\cdot W^*$ for a common $W^*$,
i.e.\ residue class only rescales the candidate limit, it does not
change its shape.
Verified on $5000$ sampled roots across four headroom levels, with
quantile ratios matching this prediction to within $2$--$9\%$; exact
only in the idealized i.i.d.\ model, as in Remark~\ref{rem:transfer-basis}.
The same relation survives unchanged when $2$ is not a primitive root
mod $q$ (the extra sterile classes of \S\ref{sec:setup}): restricted
to the surviving types $\langle2\rangle$, one can still form a
multitype pressure matrix, indexed by residue class, whose entries are
one-step transition weights rather than the deterministic map that
Example~\ref{ex:naive-fails} shows cannot exist; it is well defined
because the digit $k=\lfloor u/q\rfloor\bmod q$ that example
identifies as the missing information is itself uniform over all $q$
residues independently of the parent's class $i$, so the transition
weight out of class $i$ does not depend on $i$. This pressure matrix
is rank $1$, and its Perron eigenvalue collapses, by a
$d$-independent cancellation ($a_0$ a bijection
$\langle2\rangle\to\{1,\ldots,d\}$), to exactly $q^{\alpha-1}/(2^\alpha-1)$, the
same pressure equation~\eqref{eq:pressure-closed-form}, so both the exponent and the
scale-family direction survive intact; verified for $q=7,15$ to
$1$--$4\%$. (Whether $2$ is a primitive root mod $q$ is, for prime
$q$, governed by Artin's primitive root conjecture, open in general
but proved conditionally on GRH by \cite{Hooley1967}.)}
exactly as it
predicts it for $W_q$ in Conjecture~\ref{conj:tail-index}; the two
conjectures are expected to agree numerically at every $q$, but neither
is known to imply the other. $W$ is indexed by an actual integer root,
$W_q$ by a Haar-random $q$-adic point, and no coupling between the two
populations has been established.

At $q=3$ the evidence for Conjecture~\ref{conj:real-tree-tail} is real
and threefold: the Hill estimator and extreme-value theory on block
maxima, applied to real reverse-tree samples for the classical $q=3$
Collatz map, give $\alpha\approx2$ independently of the pressure
argument, and the direct enumeration of \S\ref{sec:empirical}
(Empirical Result~\ref{emp:real-trees}) adds a third, related line of
support (that enumeration targets
Conjecture~\ref{conj:transition-arithmetic}, the growth exponent, a
distinct claim from Conjecture~\ref{conj:real-tree-tail}, tested
further below). For $q\ge5$ the evidence for
Conjecture~\ref{conj:real-tree-tail} is markedly weaker: an early
Hill-estimator measurement of finite-depth root scale factors for
$q=5$ (600
roots, agreement to two decimal places with the predicted
$1.536\ldots$) is not statistically confirmatory at that sample size:
the estimator's true standard error is $\approx0.45$, an order of
magnitude larger than the agreement being claimed. Two further tests
bear on
Conjecture~\ref{conj:real-tree-tail} and Conjecture~\ref{conj:tail-index},
one on each.

First, on Conjecture~\ref{conj:real-tree-tail}: a larger statistical
test (5000 roots, four
headroom levels up to $10^8$) applied a battery of four estimators (a
rank-size regression with the Gabaix--Ibragimov finite-sample
correction; a bias-corrected Hill estimator following Huisman et al.;
a generalized-Pareto threshold-stability scan; and a
Clauset--Shalizi--Newman fit with a Vuong test against a truncated
log-normal alternative). The outcome is mixed:
the bias-corrected Hill estimator is stable across headroom
levels and lands close to the predicted $1.536$, but the
threshold-stability scan shows no clean plateau near the predicted
shape parameter, and the Vuong test favors the log-normal alternative
over the power law, at conventional significance, in three of the four
headroom levels tested. That stability is along one axis only, a
single estimator held fixed while headroom varies; reading the four
different estimators together by the depth of tail they each
summarize, a second and different axis, shows the apparent local tail
index
rising smoothly with depth (from $\approx1.3$ at the widest window to
$\approx2.2$ at the narrowest) without settling near any single
value: consistent with a slowly-converging pre-asymptotic regime in
which no estimator at this sample size can be treated as decisive.

Second, on Conjecture~\ref{conj:tail-index} itself: we tested it by an
exact, non-statistical route:
the same depth-$k$ recursion $Z_k(\theta;u)$ used to verify the
frozen-phase growth law above (with $\theta=\alpha_-(5)$) was evaluated
over the \emph{entire} population of residues $u\bmod 5^k$ (not a
sample), for $k$ up to $11$ ($5^{11}\approx4.9\times10^7$ residues; a
sanity check, that the population mean of $Z_k(\theta;u)$ equals
exactly $1$ for every $k$, forced by Theorem~\ref{thm:pressure} at
$\alpha=\theta$, passed to floating-point precision, confirming the
implementation). If $W_q=\lim_k Z_k(\theta;U)$ has tail index
$\iota$, the population moment $M_k(p)=\mathrm{mean}_u\,Z_k(\theta;u)^p$
should saturate in $k$ for $p<\iota$ and diverge for $p>\iota$.
The data show saturating behavior throughout $p\le1.6$ and diverging
behavior for $p\ge1.7$, which would place $\iota$ higher than the
predicted $1.536$ if taken at face value; the ratio of
successive increments has not yet stabilized for $p\le1.6$ at $k=11$
(unlike $p\ge1.7$, where it has), the classic signature of a system
still relaxing toward its asymptotic exponent. Given the observed slow
convergence rate for $q=5$ ($k^{-0.222}$ decay in the moment test, not
an effect of the linear transfer operator's spectrum, which we
verified to be exactly $\{\Lambda,0\}$ at every truncation level, i.e.\
a perfect
spectral gap with no isolated subdominant eigenvalue; the transient's
origin lies in a separate, nonlinear critical layer and remains open),
the transient decays by only about $7\%$ between
$k=8$ and $k=11$; halving it would require $k\approx250$, far beyond
what exhaustive enumeration over $5^k$ residues can reach. We therefore
cannot distinguish, from this data, between $\theta=1.536$ lying just
below a nearby true index and lying further below one. The test is
inconclusive, not disconfirming, and enumeration alone will not resolve
it.

Neither route confirms its target for $q\ge5$
(Conjecture~\ref{conj:real-tree-tail} for the first,
Conjecture~\ref{conj:tail-index} for the second), and neither refutes
it; both are more informative than the original under-powered
measurement, and both point to the same conclusion: this is a harder
statistical question than the $q=3$ case, which has stronger
finite-level support from the independent Hill/EVT measurements noted
above. Two follow-ups address the two open questions separately. The
first targets the unresolved transient in the exact test of
Conjecture~\ref{conj:tail-index}: a log-periodic
fit (motivated by the branch weights being powers of $2$) was tested
against a period derived analytically, not fitted: the relevant
renewal-theoretic dichotomy is non-arithmetic (the shift $\log_2 5$ is
irrational), so no asymptotic log-periodicity is predicted, and a
periodogram test found no power at the predicted frequencies. An
analytic handle on the linear transfer operator's spectrum, once
obtained, likewise turned out not to bear on this question, since the
transient lives in a separate nonlinear layer (the operator has a
provably perfect spectral gap). The second follow-up extends the root
battery behind Conjecture~\ref{conj:real-tree-tail}: a comparison on a
$20\times$ larger sample ($10^5$ roots)
produced, for the first time, a clean GPD threshold-stability plateau
at the predicted value, a tightly concentrated Hill estimate, and a
Vuong test that no longer favors the lognormal alternative; both were
checked against two synthetic calibration tests. This is now the
strongest evidence to date in favor of
Conjecture~\ref{conj:real-tree-tail}
for $q=5$, though not a proof of it: Vuong non-rejection is not
confirmation, and $W_v(H)$'s own median drifts from $2.04$ at
$H=10^5$ to $0.94$ at $H=10^8$, direct evidence the scale factor has
not converged at the headroom reached.

The mean drift $\log q - 2\log 2$ changes sign exactly at the
threshold $q=5$ separating the two cases of
Conjecture~\ref{conj:transition-arithmetic}: negative for $q=3$ (orbits
contract on average), positive for $q\ge 5$ (orbits expand on average).
The same sign condition governs the freezing computation of
\S\ref{sec:freezing}: at $\alpha=1$ the entropy is exactly
$s(1)=\log(4/q)$, so the same sign that decides contraction versus
expansion of orbits also decides whether the trivial root $\alpha=1$
is frozen.

\section{The tail-exponent regime at \texorpdfstring{$q=3$}{q=3}}
\label{sec:tail-regime}

At $q=3$, the second root $\alpha^\ast=2$ is numerically the boundary
value that the classical multivariate smoothing-transform theory
\cite{KoleskoMentemeier2015} would classify as the \emph{critical
case}, characterized by $m'(\alpha)=0$ at a root of $m(\alpha)=1$ (a
boundary/degenerate case in the sense of Biggins--Kyprianou). Writing
$m(\alpha):=q^{\alpha-1}/(2^\alpha-1)$ (the pressure function of
\S\ref{sec:pressure} at $q=3$), we compute
\[
  m'(1) = \ln 3 - \ln 2\cdot\frac{2^1}{2^1-1}
    \quad\text{explicitly } \approx -0.288 < 0, \qquad
  m'(2) = \ln 3 - \tfrac{4}{3}\ln 2 \approx +0.174 > 0.
\]
By Lemma~\ref{lem:log-convex}, $\log m$ is strictly convex, so the two
distinct roots of $m(\alpha)=1$ necessarily straddle its unique
minimum, so \emph{neither} can be a critical (double) root;
criticality would require the two roots to merge, which does not
happen at $q=3$. The correct classification
of $\alpha^\ast=2$ is therefore as the \emph{second root with
$m'>0$}: a Goldie/Guivarc'h-type regularly-varying tail regime
\cite{Goldie1991,BuraczewskiDamekMikosch2016}, not the critical case.
This distinction matters operationally: the critical-case uniqueness
theorem of \cite{KoleskoMentemeier2015} \emph{presupposes} exactly the
independence between subtree copies that the endogeny-barrier companion
paper identifies as the arithmetic gap obstructing a full proof
\cite{BarrierCompanion}. Even setting aside the regime mismatch,
it could not have supplied a new tool for closing that gap. Tying this
to \S\ref{sec:freezing}: ``second root with $m'>0$'' is exactly the
freezing condition of Proposition~\ref{prop:always-frozen}: the
Goldie/Guivarc'h regime name correctly classifies the algebraic root
$\alpha^\ast=2$ of $m(\alpha)=1$, but by itself says nothing about
whether this root governs the \emph{quenched} tail of our specific
arithmetic recursion. That is exactly the question
Proposition~\ref{prop:always-frozen}, Conjecture~\ref{conj:tail-index},
and Conjecture~\ref{conj:real-tree-tail} address directly.

\section{Empirical confirmation on real reverse trees}
\label{sec:empirical}

\begin{empirical}[Direct enumeration on real reverse trees]\label{emp:real-trees}
Independent depth-first enumeration of the reverse
trees for $q=5$ and $q=7$, following the exact admissibility rule of
\S\ref{sec:setup}, provides direct empirical support for
Conjecture~\ref{conj:transition-arithmetic}:
counting slopes per decade of $3$ independently sampled roots gave
$\approx 0.62$--$0.66$ for $q=5$ (predicted $0.650919$) and
$\approx 0.36$--$0.39$ for $q=7$ (predicted $0.373501$).
\end{empirical}

A more careful,
statistically calibrated measurement for $q=5$, separating this
prediction from a competing one, is the subject of a companion paper
\cite{KLVolkovCompanion}. As already noted above for $q=3$ (where
$w_2(1)=1$ is a fixed point), the smallest-branch map
$u\mapsto w_{a_0(u\bmod q)}(u)$ has finite cycles for $q=5$ (e.g.\
$1\mapsto3\mapsto1$, since $w_4(1)=3$ and $w_1(3)=1$) and $q=7$; naive
enumeration would loop on nodes that reach one. Our implementation
excludes all such cycle members explicitly (precomputed for each $q$)
and deduplicates visited nodes, instead of only sampling
roots far from them.

\begin{remark}[Bibliographic calibration]\label{rem:novelty-109}
Neither the numerical value $\alpha^\ast(5)=0.650919$ nor the
qualitative transition at $q\ge5$ is new, and a later result rigorously
excludes $q\ge5$ in the forward direction for the analogous $p=2$
generalization \cite{GoncalvesGreenfeldMadrid2022}. Wirsching
\cite[Ch.~III, \S5, following Theorem~5.2]{Wirsching1998} anticipated
the same qualitative transition at $p>4$ heuristically in 1998, in his
own terms (a sign flip in the index of the maximal term $f_p(n)$,
leading him to expect vanishing asymptotic density for the weak
components of the Collatz graph of $T_p$),
and Kontorovich--Lagarias \cite{KontorovichLagarias2010} had already
computed the same quantity
(denoted $\eta_{5,BP}$ in their Theorem~8.10) via a different route,
large-deviations methods for a branching random walk, that is
numerically identical to our pressure-equation derivation. Given
\S\ref{sec:freezing}, this difference in route matters beyond giving
an alternative derivation: a large-deviations argument for a
branching random walk is, in the appropriate sense, a quenched
statement, whereas our pressure identity (Theorem~\ref{thm:pressure})
is annealed. That the two agree numerically at $\alpha_-(5)$ is
consistent with $\alpha_-(5)$ being unfrozen
(Proposition~\ref{prop:always-frozen}), where quenched and annealed
are expected to coincide; it does not by itself certify that our
annealed route would agree with a quenched one at a frozen root. The
closed-form equation~\eqref{eq:pressure-eq} is new: it unifies the
family for arbitrary $q$ in a single algebraic
identity, with a rigorous annealed derivation and
independent numerical confirmation.
\end{remark}

\section*{Data Availability Statement}

Code and data reproducing every claim in this paper, with a README
mapping each script to the section and result it verifies, are at
\url{https://github.com/faculdade/collatz-qx1-pressure}.

\section*{Acknowledgments}

AI assistance (Claude, Anthropic) was used for textual review and
translation.

\bibliographystyle{plain}
\begin{thebibliography}{99}

\bibitem{BarrierCompanion} R.\ A.\ Tavares, \emph{The $qx+1$
  generalization of Tao's Syracuse measure: an endogeny barrier},
  companion paper, in preparation.

\bibitem{KLVolkovCompanion} R.\ A.\ Tavares, \emph{A calibrated
  separation of the Kontorovich--Lagarias and Volkov exponent
  predictions for the $5x+1$ reverse tree}, companion paper, in
  preparation.

\bibitem{KontorovichLagarias2010} A.\ V.\ Kontorovich and J.\ C.\
  Lagarias, \emph{Stochastic models for the $3x+1$ and $5x+1$
  problems}, in: \emph{The Ultimate Challenge: the $3x+1$ Problem}
  (J.\ C.\ Lagarias, ed.), American Mathematical Society, 2010.

\bibitem{ApplegateLagarias1995I} D.\ Applegate and J.\ C.\ Lagarias,
  \emph{Density bounds for the $3x+1$ problem I: tree-search method},
  Mathematics of Computation \textbf{64} (1995), 411--426.

\bibitem{ApplegateLagarias1995II} D.\ Applegate and J.\ C.\ Lagarias,
  \emph{Density bounds for the $3x+1$ problem II: Krasikov
  inequalities}, Mathematics of Computation \textbf{64} (1995),
  427--438.

\bibitem{KrasikovLagarias2003} I.\ Krasikov and J.\ C.\ Lagarias,
  \emph{Bounds for the $3x+1$ problem using difference inequalities},
  Acta Arithmetica \textbf{109}(3) (2003), 237--258.

\bibitem{Wirsching1998} G.\ J.\ Wirsching, \emph{The Dynamical System
  Generated by the $3n+1$ Function}, Lecture Notes in Mathematics
  \textbf{1681}, Springer, 1998.

\bibitem{GoncalvesGreenfeldMadrid2022} F.\ Gon\c{c}alves, R.\
  Greenfeld, and J.\ Madrid, \emph{Generalized Collatz maps with almost
  bounded orbits}, arXiv:2111.06170, 2022.

\bibitem{Biggins1992} J.\ D.\ Biggins, \emph{Uniform convergence of
  martingales in the branching random walk}, Annals of Probability
  \textbf{20}(1) (1992), 137--151.

\bibitem{Nerman1981} O.\ Nerman, \emph{On the convergence of
  supercritical general (C-M-J) branching processes}, Zeitschrift
  f\"ur Wahrscheinlichkeitstheorie und verwandte Gebiete
  \textbf{57}(3) (1981), 365--395.

\bibitem{Goldie1991} C.\ M.\ Goldie, \emph{Implicit renewal theory and
  tails of solutions of random equations}, Annals of Applied
  Probability \textbf{1}(1) (1991), 126--166.

\bibitem{JelenkovicOlveraCravioto2012} P.\ R.\ Jelenkovi\'c and
  M.\ Olvera-Cravioto, \emph{Implicit renewal theorem for trees with
  general weights}, Stochastic Processes and their Applications
  \textbf{122}(9) (2012), 3209--3238.

\bibitem{AlsmeyerBigginsMeiners2012} G.\ Alsmeyer, J.\ D.\ Biggins, and
  M.\ Meiners, \emph{The functional equation of the smoothing
  transform}, Annals of Probability \textbf{40}(5) (2012), 2069--2105.

\bibitem{VillemonaisZalduendo2025} D.\ Villemonais and N.\ Zalduendo,
  \emph{Convergence of weighted branching processes}, arXiv:2512.07653
  (2025).

\bibitem{Liu2000} Q.\ Liu, \emph{On generalized multiplicative
  cascades}, Stochastic Processes and their Applications
  \textbf{86}(2) (2000), 263--286.

\bibitem{Hooley1967} C.\ Hooley, \emph{On Artin's conjecture}, Journal
  f\"ur die reine und angewandte Mathematik \textbf{225} (1967),
  209--220.

\bibitem{KoleskoMentemeier2015} K.\ Kolesko and S.\ Mentemeier,
  \emph{Fixed points of the multivariate smoothing transform: the
  critical case}, Electronic Journal of Probability \textbf{20}
  (2015), \href{https://doi.org/10.1214/EJP.v20-4022}{doi:10.1214/EJP.v20-4022}.

\bibitem{BuraczewskiDamekMikosch2016} D.\ Buraczewski, E.\ Damek, and
  T.\ Mikosch, \emph{Stochastic Models with Power-Law Tails: The
  Equation $X=AX+B$}, Springer, 2016.

\end{thebibliography}

\end{document}
